% ================= CHƯƠNG 4: KẾT QUẢ VÀ ĐÁNH GIÁ =================
\chapter{KẾT QUẢ VÀ ĐÁNH GIÁ}
\label{chap:ket_qua_danh_gia}

\section{Kết quả và đánh giá về chức năng chung}
\label{sec:danh_gia_chuc_nang_chung}

Qua quá trình thử nghiệm nghiêm ngặt trên cả môi trường máy ảo Android Emulator và thiết bị Android vật lý thực tế (Samsung Galaxy S21, Xiaomi Redmi Note 12), các phân hệ chức năng nền tảng của ứng dụng đều vận hành ổn định và đạt kết quả đề ra:

\begin{enumerate}
    \item \textbf{Phân hệ Xác thực và Đăng nhập:}
    \begin{itemize}
        \item Đăng nhập bằng tài khoản Email/Mật khẩu diễn ra chuẩn xác; thời gian xác thực trung bình với máy chủ Firebase Authentication dao động từ 1.0 đến 1.4 giây tùy điều kiện mạng.
        \item Tính năng đăng nhập một chạm với Google thông qua \texttt{CredentialManager} hoạt động mượt mà, tự động lấy tên hiển thị và ảnh đại diện để điền sẵn vào hồ sơ người dùng mới.
        \item Luồng khôi phục mật khẩu thông qua mã xác minh OTP hoạt động tin cậy, tự động giới hạn thời gian hiệu lực của mã trong 5 phút để bảo đảm an toàn thông tin.
    \end{itemize}

    \item \textbf{Phân hệ Trang chủ và Bảng điều khiển (Home Dashboard):}
    \begin{itemize}
        \item Màn hình chính tổng hợp trực quan tình trạng thể trạng của thú cưng đang được chọn: Cân nặng hiện tại, thân nhiệt ($38.4^\circ\text{C}$), mức độ vận động hàng ngày ($8.240$ bước) và chế độ ăn uống ($2/2$ bữa đạt yêu cầu).
        \item Danh sách công việc cần làm trong ngày được hiển thị dưới dạng thẻ ngắn gọn, cho phép người dùng tick chọn hoàn thành ngay trên màn hình chính mà không cần chuyển màn hình.
    \end{itemize}
\end{enumerate}

\section{Kết quả và đánh giá về các chức năng chuyên sâu}
\label{sec:danh_gia_chuc_nang_chuyen_sau}

\subsection{Chức năng quản lý hồ sơ thú cưng}
Hệ thống cho phép người dùng khởi tạo và quản lý không giới hạn số lượng thú cưng trên một tài khoản. Việc chuyển đổi giữa các thú cưng thông qua dải biểu tượng tròn ở đầu màn hình diễn ra tức thì (dưới 50ms) nhờ việc lưu trữ trạng thái \texttt{currentSelectedPetId} trong \texttt{PetRepository}. Khi chuyển thú cưng, toàn bộ dữ liệu tiêm chủng, đơn thuốc và lịch trình tương ứng đều được lọc lại một cách tự động.

\subsection{Sổ sức khỏe điện tử (Health Hub)}
Đây là phân hệ cốt lõi giải quyết bài toán thất lạc dữ liệu y tế của người nuôi thú cưng:
\begin{itemize}
    \item \textbf{Tiêm chủng (Vaccination):} Quản lý rõ ràng danh sách các mũi tiêm với nhãn phân biệt màu sắc: màu xanh lá cho trạng thái ``Đã tiêm'' và màu cam nổi bật cảnh báo cho các mũi ``Nhắc lại''. Người dùng dễ dàng theo dõi thời hạn để đưa thú cưng đi tiêm chủng đúng kỳ.
    \item \textbf{Đơn thuốc (Medication):} Cung cấp thông tin chi tiết tên thuốc, liều lượng dùng trong ngày và đếm lùi số lượng viên còn lại, giúp người nuôi không bị quên hoặc nhầm lẫn liều lượng.
    \item \textbf{Nhật ký bệnh án (Medical Record Book):} Lưu trữ các đợt khám bệnh bất thường cùng hình ảnh chứng minh y khoa, tạo thành một cuốn bệnh án điện tử có cấu trúc khoa học phục vụ đắc lực cho bác sĩ thú y khi tiếp nhận khám lâm sàng.
    \item \textbf{Lịch sử cân nặng và Khám tổng quát:} Ghi chép chi tiết chỉ số cân nặng qua các tháng, hỗ trợ phát hiện sớm các dấu hiệu suy dinh dưỡng hoặc béo phì ở vật nuôi.
\end{itemize}

\subsection{Chức năng Lịch trình thông minh (Calendar \& Reminders)}
Lịch trình hiển thị theo cấu trúc dạng tuần thông minh. Các sự kiện chăm sóc được phân loại rõ ràng theo từng loại biểu tượng (khám bệnh, tiêm ngừa, uống thuốc, tắm rửa, giờ ăn). Thuật toán sắp xếp tự động đẩy các sự kiện chưa hoàn thành lên vị trí ưu tiên để thu hút sự chú ý của người nuôi.

\subsection{Bản đồ dịch vụ thú y (Google Maps Integration)}
\begin{itemize}
    \item Tích hợp bản đồ vệ tinh và bản đồ giao thông tương tác cao, tự động định vị vị trí người dùng thông qua GPS của thiết bị.
    \item Hệ thống lọc theo 3 nhóm danh mục chính (Phòng khám Thú y, Spa chăm sóc, Cửa hàng thức ăn).
    \item Thuật toán Haversine tính toán khoảng cách chuẩn xác, hiển thị ngay trên thẻ thông tin tóm tắt kèm nút gọi điện khẩn cấp và nút mở ứng dụng Google Maps chỉ đường.
\end{itemize}

\begin{table}[H]
    \caption{Bảng đánh giá mức độ hoàn thiện các phân hệ chức năng thực tế}
    \label{tab:danh_gia_chuc_nang}
    \centering
    \renewcommand{\arraystretch}{1.2}
    \begin{tabularx}{\textwidth}{|l|l|c|X|}
        \hline
        \textbf{Phân hệ} & \textbf{Chức năng kiểm thử} & \textbf{Tỷ lệ đạt} & \textbf{Đánh giá kết quả} \\ \hline
        \multirow{3}{*}{Xác thực} 
        & Đăng ký / Đăng nhập Email & 100\% & Hoạt động ổn định, kiểm tra hợp lệ tốt. \\ \cline{2-4}
        & Google Sign-In (Credential Manager) & 100\% & Nhanh chóng, tự động lấy thông tin hồ sơ. \\ \cline{2-4}
        & Quên mật khẩu / Xác minh OTP & 100\% & Mã OTP gửi nhanh, kiểm tra thời hạn chuẩn. \\ \hline
        
        \multirow{2}{*}{Thú cưng} 
        & Thêm thú cưng thủ công & 100\% & Đầy đủ trường dữ liệu, lưu trữ thành công. \\ \cline{2-4}
        & Chuyển đổi và chỉnh sửa thú cưng & 100\% & Cập nhật tức thời lên giao diện chính. \\ \hline

        \multirow{4}{*}{Sức khỏe} 
        & Quản lý sổ tiêm chủng & 100\% & Phân loại rõ ràng mũi đã tiêm và nhắc lại. \\ \cline{2-4}
        & Quản lý đơn thuốc điều trị & 100\% & Hiển thị liều lượng và giờ uống chính xác. \\ \cline{2-4}
        & Nhật ký bệnh án & 100\% & Lưu trữ đính kèm hình ảnh và cảm xúc. \\ \cline{2-4}
        & Theo dõi cân nặng & 100\% & Biểu diễn số liệu thể trạng trực quan. \\ \hline

        \multirow{2}{*}{Lịch trình} 
        & Thêm mới và quản lý sự kiện & 100\% & Phân loại đa dạng theo hoạt động chăm sóc. \\ \cline{2-4}
        & Đánh dấu hoàn thành việc cần làm & 100\% & Cập nhật thanh tiến độ trong ngày chính xác. \\ \hline

        \multirow{2}{*}{Bản đồ} 
        & Định vị GPS và hiển thị Marker & 100\% & Tọa độ chính xác, bản đồ tương tác mượt. \\ \cline{2-4}
        & Lọc địa điểm và tính khoảng cách & 100\% & Công thức Haversine tính khoảng cách chuẩn. \\ \hline
    \end{tabularx}
    \source{Nhóm tác giả tổng hợp từ kết quả chạy kiểm thử thực tế}
\end{table}

\section{Đánh giá các chỉ tiêu phi chức năng}
\label{sec:danh_gia_phi_chuc_nang}

\subsection{Độ mượt mà và Tài nguyên phần cứng}
Nhóm sử dụng công cụ Android Profiler trong Android Studio để đo đạc các thông số vận hành thực tế trên thiết bị Samsung Galaxy S21:
\begin{itemize}
    \item \textbf{Tốc độ khung hình (Frame Rate):} Tốc độ dựng khung hình duy trì ổn định ở mức từ 58.5 đến 60 FPS trong suốt quá trình cuộn danh sách (LazyColumn) và chuyển tab, không ghi nhận khung hình bị đóng băng (frozen frame).
    \item \textbf{Mức độ tiêu thụ RAM:} Ứng dụng chiếm dụng trung bình từ 85MB đến 130MB RAM khi chạy các tác vụ cơ bản và đạt tối đa khoảng 175MB khi mở bản đồ Google Maps có tải nhiều hình ảnh bitmap, hoàn toàn nằm trong ngưỡng an toàn của hệ điều hành.
    \item \textbf{Mức độ tiêu hao CPU:} CPU duy trì dưới 5\% ở trạng thái chờ và từ 12\% đến 22\% khi người dùng thao tác chuyển đổi trang liên tục hoặc tương tác bản đồ.
\end{itemize}

\subsection{Tính sẵn sàng và Khả năng chịu lỗi ngoại lệ}
Khi thử nghiệm ngắt kết nối Internet trong lúc thực hiện thao tác thêm mũi tiêm hoặc chỉnh sửa lịch trình:
\begin{itemize}
    \item Ứng dụng không phát sinh lỗi văng ứng dụng (Crash).
    \item Dữ liệu vẫn được hiển thị và cập nhật bình thường trên giao diện nhờ tầng nhớ đệm ngoại tuyến của Firestore SDK.
    \item Khi tái kết nối Internet, toàn bộ dữ liệu mới được đồng bộ lên máy chủ đám mây trong vòng chưa đầy 1 giây mà không yêu cầu người dùng phải thao tác lại.
\end{itemize}

\subsection{Trải nghiệm người dùng và Chuẩn thiết kế}
Ứng dụng tuân thủ nghiêm ngặt bảng hướng dẫn Material Design 3. Độ tương phản màu sắc của chữ viết trên nền thẻ đạt tỷ lệ tương phản tối thiểu 4.5:1, kích thước các vùng chạm (Touch Target) của nút bấm đều đạt chuẩn tối thiểu $48\times 48\text{ dp}$, giúp người dùng dễ dàng thao tác một tay trên các thiết bị di động màn hình lớn.
